\subsection{Das Messaging-Konzept}\label{sec:runtime_messaging}

Der Nachrichtenaustausch zwischen Systemen und Laufzeitkomponenten erfolgt über den \texttt{UpdateContext}, der als Kontextobjekt Schnittstellen zum Lesen und Ablegen von Nachrichten bereitstellt (siehe Abschnitt~\ref{sec:runtime_gameloop}).
Darunter befindet sich der Zugriff auf einen zentralen und Engine-spezifischen \texttt{CommandBuffer}, der - basierend auf zur Kompilierzeit definierten Befehlstypen - relevante \textit{Commands} zur Ablage und Weiterverarbeitung bereitstellt.\footnote{
    Die Angabe der zur Kompilierzeit bekannten Befehlstypen ermöglicht einen Performance-Vorteil, der mit erhöhtem Bootstrapping- und Wartungsaufwand bei der Einführung neuer Commands einhergeht, wie in Abschnitt~\ref{sec:meilenstein_5} beschrieben.
}

Die Command-Verarbeitung stützt sich auf die \texttt{CommandHandlerRegistry}:
Bei der Abarbeitung (\textit{Flush)} des \texttt{TypedCommandBuffer} wird für jeden Befehlstyp geprüft, ob ein registrierter Handler vorliegt.
Ist dies der Fall, wird der Command an diesen delegiert; andernfalls wird - sofern der Command das \texttt{ExecutableCommand}-Concept erfüllt - dessen \texttt{execute}-Methode direkt aufgerufen.\\

Die Registrierung eines Handlers erfolgt über die \texttt{GameWorld}:

\begin{verbatim}
gameWorld.registerCommandHandler<SpawnCommand, DespawnCommand>(handler);
\end{verbatim}

Die \texttt{GameWorld} nutzt hierbei \textit{Variadic Templates} zusammen mit \textit{Fold Expressions}~\cite[S. 82-118]{Ban22}, um für jeden übergebenen Command-Typ einen zugehörigen Handler-Eintrag in der \texttt{ResourceRegistry} der \texttt{GameWorld} anzulegen.
Auch hier werden die Handler erneut über \textit{Concepts} (\texttt{IsCommandHandlerLike}) validiert, um sicherzustellen, dass für jeden angegebenen Typ eine passende \texttt{submit}-Funktion implementiert ist.
Intern speichert die Registry eine \textit{type erased} Indirektion, sodass Methoden des Handlers ohne Vererbung von einer gemeinsamen Basisklasse aufgerufen werden können.

\vspace{5mm}
\begin{tcolorbox}[title={Verarbeitungsreihenfolge von Befehlen}, colback=black!5, colframe=black]
    In der \texttt{GameWorld} registrierte \texttt{CommandBuffer} werden von der \texttt{GameLoop} bei jedem Phasenwechsel geflusht.
    Darüber hinaus ist es möglich, über \texttt{CommitPoints} ein Flushen innerhalb einer Phase zwischen zwei Pässen auszulösen.\\
    Dabei wird der Write-Buffer des \texttt{CommandBuffer} mit dem Read-Buffer getauscht, was für nachfolgende Pässe, die über den \texttt{UpdateContext} Commands eines bestimmten Typs erwarten, relevant ist:
    Systeme müssen daher so angeordnet werden, dass die Reihenfolge der Verarbeitung die erfolgten Buffer-Swaps berücksichtigt.
\end{tcolorbox}
\vspace{5mm}

Neben Commands steht mit den \textit{Events} ein zweiter Kommunikationskanal zur Verfügung.
Die Unterteilung in Commands und Events ist dabei nicht strikt, sondern eher konzeptionell:
Commands sind immer dann vorgesehen, wenn der Weltzustand deterministisch verändert werden soll.
Events hingegen sind Informationen, die eine Weiterverarbeitung von Zustandsänderungen ermöglichen sollen, beispielsweise für visuelle oder auditive Effekte, und damit Fakten beziehungsweise Signale zwischen Systemen transportieren.\\

Der \texttt{UpdateContext} bietet hierbei Zugriff auf drei Ebenen von Event-Bussen, die ebenfalls als Double Buffer realisiert sind (vgl.\ Abschnitt~\ref{sec:runtime_gameloop}):

\begin{itemize}
    \itemsep0.3em
    \item \textbf{Pass-Events} (\texttt{pushPass}/\texttt{readPass}): Werden an \texttt{CommitPoints} innerhalb einer Phase getauscht und stehen nachfolgenden Pässen derselben Phase zur Verfügung.
    \item \textbf{Phase-Events} (\texttt{pushPhase}/\texttt{readPhase}): Werden an Phasengrenzen getauscht und sind in der jeweils nächsten Phase lesbar.
    \item \textbf{Frame-Events} (\texttt{pushFrame}/\texttt{readFrame}): Werden am Ende eines Frames getauscht und erlauben eine frameübergreifende Kommunikation.
\end{itemize}

Durch die Aufteilung in diese drei Ebenen und eine entsprechende Strukturierung der GameLoop lässt sich die Reichweite einer Nachricht gezielt eingrenzen.